\section{Formalización del entorno: POMDP, espacio de estados y espacio de acciones}

\subsection{Por qué formalizar}
El expositor modela al Web Agent como un POMDP, un paso muy importante. Esto nos permite describir el problema con un lenguaje unificado de aprendizaje por refuerzo y búsqueda, en lugar de tratar el comportamiento del sistema como un conjunto de scripts heurísticos.

\begin{figure}[H]
\centering
\includegraphics[width=0.82\textwidth]{frame-02-pomdp.jpg}
\caption{Video aproximadamente en 00:13:22: el expositor escribe el problema de interacción con la página web como states / actions / observations / transitions}
\end{figure}

\begin{knowledgebox}{Explicación intuitiva del POMDP en este contexto}
\begin{itemize}
    \item \textbf{State}: el estado real de la página web (incluida la información no visible);
    \item \textbf{Observation}: el contenido actualmente visible para el Agent (captura de pantalla, DOM, árbol de accesibilidad, etc.);
    \item \textbf{Action}: clic, escritura, desplazamiento, cambio de pestaña, detención, etc.;
    \item \textbf{Reward}: si la tarea se completó correctamente; suele ser escasa.
\end{itemize}
\end{knowledgebox}

\subsection{El diseño del espacio de acciones determina el límite superior del sistema}
Si el espacio de acciones es demasiado grueso, la capacidad expresiva resulta insuficiente; si es demasiado fino, el espacio de búsqueda explota. En la práctica suele necesitarse un diseño jerárquico de «acciones semánticas + acciones de bajo nivel»; por ejemplo, un subobjetivo como «buscar restaurante» lo entrega el planner, y el executor lo descompone en clics y escritura.

\begin{importantbox}{Principios de diseño del espacio de acciones}
\begin{enumerate}
    \item Las acciones atómicas deben ser ejecutables y verificables;
    \item Las acciones de alto nivel deben reducir la profundidad de búsqueda;
    \item La acción de detención debe ir acompañada de un criterio confiable para determinar la finalización;
    \item Cada acción debe llevar un registro rastreable, para facilitar la reproducción y la depuración.
\end{enumerate}
\end{importantbox}

\begin{table}[H]
\centering
\small
\begin{tabular}{p{3.2cm}p{4.1cm}p{3.7cm}}
\toprule
\textbf{Nivel de acción} & \textbf{Acciones típicas} & \textbf{Riesgo principal} \\
\midrule
Acciones de bajo nivel & click / type / scroll / hover & pasos largos; los errores locales se acumulan fácilmente \\
Acciones de nivel medio & open-site / query / filter & ambigüedad semántica; dificultad para vincular parámetros \\
Acciones de alto nivel & complete-subtask / stop & detención prematura o juicio erróneo de finalización \\
\bottomrule
\end{tabular}
\caption{Niveles de acción del Web Agent y sus riesgos}
\end{table}

\subsection{Retos de entrenamiento y evaluación bajo recompensas escasas}
En las tareas web, muchos reward solo se otorgan al final, por lo que la asignación de crédito es muy difícil. Si el sistema depende únicamente de la señal final, difícilmente aprende la capacidad de corrección fina en los pasos intermedios.

\begin{warningbox}{Efectos secundarios de la recompensa escasa}
\begin{itemize}
    \item El entrenamiento es inestable, con alta varianza en la actualización de la política;
    \item El modelo tiende a aprender shortcuts oportunistas;
    \item En las tareas de cadena larga, los errores intermedios son difíciles de penalizar a tiempo.
\end{itemize}
\end{warningbox}

\begin{figure}[H]
\centering
\includegraphics[width=0.82\textwidth]{frame-03-trajectory.jpg}
\caption{Video aproximadamente en 00:16:03: el expositor muestra la trayectoria de una tarea compleja y enfatiza la coherencia entre el razonamiento de cadena larga y la ejecución}
\end{figure}

\subsection{Resumen de la sección}
La formalización mediante POMDP hace que la tarea del Agent pase de ser «ingeniería de scripts» a ser un «problema de decisión». Pero la recompensa escasa y la explosión del espacio de acciones hacen que, con un solo paso de inferencia hacia adelante, sea difícil completar tareas complejas de manera estable.

